\chapter{Tailbone (Coccyx) Pain}

\section*{Definition}
Tailbone pain is discomfort at the base of the spine, often worse with sitting, rising, bowel movements, or direct pressure. It commonly follows a fall or prolonged pressure but may have other causes.

\section*{When to See a Doctor}
Seek urgent care for fever, redness or drainage, severe trauma, progressive numbness or weakness, saddle anesthesia, or loss of bladder or bowel control. Arrange prompt evaluation for persistent, recurrent, unexplained, or function-limiting symptoms.

\section*{How It Develops}
The symptom results from irritation, inflammation, altered organ function, injury, infection, or disruption of normal neurologic or vascular signaling. Age, pregnancy status, onset, distribution, severity, and associated findings determine the urgency and likely causes.

\section*{Causes}
\begin{itemize}
  \item Bruising or fracture, childbirth, prolonged sitting, repetitive strain, coccyx instability, pilonidal disease, arthritis, infection, and rarely tumor.
  \item Medication effects, environmental exposures, and underlying chronic disease may also contribute.
  \item Serious causes are less common but must be considered when symptoms are sudden, progressive, severe, or associated with systemic illness.
\end{itemize}

\section*{Related Symptoms}
\begin{itemize}
  \item Pain, fever, fatigue, nausea, weakness, or changes in normal function
  \item Bleeding, swelling, discharge, breathing difficulty, or neurologic symptoms when a serious cause is present
\end{itemize}

\section*{Medical Evaluation}
\begin{itemize}
  \item History addresses onset, duration, severity, triggers, injuries, exposures, medications, medical conditions, age, and pregnancy status when relevant.
  \item Examination includes vital signs and focused assessment of the relevant organ system, neurologic status, hydration, and function.
  \item Testing may include blood or urine studies, cultures, imaging, physiologic testing, fetal or pediatric assessment, or specialist examination.
\end{itemize}

\section*{Treatment Options}
Treatment is directed at the cause and may include supportive care, targeted medication, treatment of infection or inflammation, correction of metabolic disease, physical therapy, or specialist procedures. Persistent symptoms require reassessment.

\section*{Self-Care and Home Management}
Avoid known triggers, protect the affected area, maintain appropriate hydration, and record timing and associated symptoms. Use only age- or pregnancy-appropriate medicines recommended by a clinician. Do not delay emergency care while trying home treatment.

\section*{References}
\begin{thebibliography}{99}
\bibitem{tailbone_pain:jaaos} Fogel GR, Cunningham PY 3rd, Esses SI. Coccygodynia: evaluation and management. \textit{J Am Acad Orthop Surg}. 2004;12:49--54.
\bibitem{tailbone_pain:pain-pract} Patijn J, Janssen M, Hayek S, et al. Coccygodynia. \textit{Pain Pract}. 2010;10:554--559.
\end{thebibliography}
